\documentclass[10pt,a4paper,twocolumn]{article}

\usepackage[utf8]{inputenc}
\usepackage[T1]{fontenc}
\usepackage{lmodern}
\usepackage[margin=1.8cm,columnsep=0.7cm]{geometry}
\usepackage{amsmath,amssymb,amsthm}
\usepackage[round,authoryear]{natbib}
\usepackage{xcolor}
\usepackage{hyperref}
\usepackage{enumitem}
\usepackage{booktabs}

\hypersetup{colorlinks=true,linkcolor=blue!60!black,
  citecolor=green!50!black,urlcolor=blue!70!black}

\theoremstyle{plain}
\newtheorem{theorem}{Theorem}
\newtheorem{proposition}[theorem]{Proposition}
\newtheorem{lemma}[theorem]{Lemma}
\newtheorem{corollary}[theorem]{Corollary}
\newtheorem{conjecture}[theorem]{Conjecture}
\theoremstyle{definition}
\newtheorem{definition}{Definition}
\newtheorem{remark}{Remark}

\newcommand{\dd}{\mathrm{d}}
\newcommand{\ee}{\mathrm{e}}
\newcommand{\eps}{\varepsilon}
\newcommand{\Cstr}[1]{\mathcal{C}_{#1}}
\newcommand{\Nalg}{\mathbb{N}\text{-algebraic}}

\usepackage{mdframed}
\newmdenv[backgroundcolor=blue!5,linecolor=blue!40,
          linewidth=1pt,innertopmargin=6pt,innerbottommargin=6pt,
          innerleftmargin=8pt,innerrightmargin=8pt]{slotbox}

\title{The $\Cstr{3}$ Stratum in Non-DP Universality\\[2pt]
  \large Slotting Manna and Conserved-DP, and what currently resists slotting}

\author{S. Amarteifio\\
  \small\textit{Percolation Labs}}
\date{April 2026}

\begin{document}
\maketitle

\begin{abstract}
The non-directed-percolation (non-DP) universality classes catalogued
by~\citet{Odor2004} --- Manna, voter, conserved-DP (C-DP), and a
handful of others --- have resisted a unified organising principle.
We apply the CFAC stratification theorem~\citep{AmarteifioCFAC2026},
which organises polynomial Dyson--Schwinger equations by the Puiseux
order $k$ of their dominant branch point via $\tau = 1 + 1/k + \gamma_k(d)$,
where the first two terms are a $d$-independent skeleton and $\gamma_k(d)$
is a $d$-dimensional dressing computable from the rank-$k$ bridge
function.  Two observations:
\textbf{(i)} The Manna class ($\tau \approx 1.29$) and conserved-DP
($\tau \approx 1.35$) both sit within $3\%$ of the $\Cstr{3}$ skeleton
$\tau = 4/3$, flanking it from opposite sides --- a non-trivial
prediction that the two classes are a common $k=3$ stratum split only
by their $\gamma_3(d)$ dressings.  Non-dyadic $k=3$ is excluded by
Banderier--Drmota for $\Nalg$ DSEs but permitted by signed/annihilation
dynamics, which both conserved classes plausibly exhibit.
\textbf{(ii)} We predict the Manna--C-DP $\tau$-gap from the rank-3
bridge function without fitting parameters and compare to published
simulations.
The paper also documents three directions --- LERW in $d=3$, TAP
complexity in spin glasses, KPZ upper-tail universality --- where we
\emph{believe} the current CFAC machinery does not reach, with stated
structural reasons.  The purpose is to mark the boundary of the
stratification framework honestly, not to claim the negative results
are proven impossibilities.
\end{abstract}

\tableofcontents

% ==================================================================
\section{Introduction}
\label{sec:intro}
% ==================================================================

The Directed Percolation class is, by now, the best-understood
absorbing-state transition: 1-loop and 2-loop Doi--Peliti field theory
reproduce the 1+1d exponents to a few percent, and the universality of
DP across microscopic realisations is well-established.
Outside DP, the picture is more fragmented.
\citet{Odor2004} catalogued a handful of additional classes:
the \emph{Manna} class of stochastic sandpiles, the \emph{voter} model
in its 1d degenerate limit, the \emph{conserved-DP} class arising
when a conserved background field is coupled to a DP order parameter,
and a few others.  Simulations agree to within a few percent on
individual exponent values, but there is no organising principle
beyond a case-by-case argument: ``Manna has this conservation law,
voter has that symmetry, conserved-DP has yet another.''

The CFAC stratification theorem~\citep{AmarteifioCFAC2026} offers one
candidate for such a principle.  Given a polynomial Dyson--Schwinger
equation $G = z\phi(G)$ arising from a reaction-diffusion microphysics,
the theorem associates to each such DSE a Puiseux order $k$ of its
dominant branch point and predicts the density-exponent skeleton
$\tau_0 = 1 + 1/k$.  The $d$-dependent correction $\gamma_k(d)$ arises
at one loop from the rank-$k$ bridge function and shifts the measured
$\tau$ by an $O(\eps)$ amount in $\eps = d_c - d$.

In this note we test the hypothesis that Manna and conserved-DP share
a common $k=3$ stratum, with the $4\%$ gap between their simulation
values $\tau_{\rm Manna} \approx 1.29$ and $\tau_{\rm C-DP} \approx 1.35$
accounted for entirely by the difference in their one-loop
$\gamma_3(d=1+1)$.  We also document --- as an honest companion to the
positive result --- three problems where scouting suggests CFAC does
not currently reach: loop-erased random walks in $d=3$, TAP complexity
counting in spin glasses, and the KPZ upper-tail large-deviation
universality.  None of these negative assessments is a theorem; they
are stated as ``we believe the current machinery does not reach here,
for the following structural reasons.''

\paragraph{Contents.}
Section~\ref{sec:recap} recalls the stratification theorem.
Section~\ref{sec:slotting} states the $\Cstr{3}$ slotting claim for
Manna and conserved-DP.
Section~\ref{sec:gamma3} computes the rank-3 bridge
$\gamma_3(d=1+1)$ and predicts the $\tau$-gap.
Section~\ref{sec:sim} compares to published simulations.
Section~\ref{sec:boundary} documents what currently resists slotting
and why.
Section~\ref{sec:discussion} concludes.

% ==================================================================
\section{The CFAC stratification theorem}
\label{sec:recap}
% ==================================================================

We quote only the structure needed here; full proofs appear
in~\citep{AmarteifioCFAC2026}.

\begin{definition}[Polynomial DSE]
A \emph{polynomial Dyson--Schwinger equation} is $G = z\phi(G)$ where
$\phi \in \mathbb{Q}[G]$ is a polynomial of degree $\ge 2$ with
non-negative coefficients (or, in extensions, with signed coefficients
subject to a stability condition).
\end{definition}

\begin{theorem}[Stratification,~\citep{AmarteifioCFAC2026}]
\label{thm:strat}
Let $G(z)$ satisfy $G = z\phi(G)$ with $\phi$ polynomial.  Then the
dominant singularity $z_\star$ is a Puiseux branch point of order
$k \in \{2, 3, 4, \ldots\}$, and the density exponent satisfies
\begin{equation}\label{eq:tau_strat}
  \tau(d) = 1 + \frac{1}{k} + \gamma_k(d), \qquad
  \gamma_k(d) = c_k\,\eps + O(\eps^2),
\end{equation}
where $\eps = d_c - d$ and $c_k$ is a rational multiple of the rank-$k$
bridge function $B_k$ evaluated at the one-loop fixed point.
\textbf{Banderier--Drmota constraint:} for $\Nalg$ DSEs (non-negative
polynomial coefficients), the dominant $k$ lies in the dyadic ladder
$\{2, 4, 8, \ldots\}$.  For signed/annihilation DSEs, non-dyadic
integer $k$ is accessible.
\end{theorem}

The skeleton term $\tau_0 = 1 + 1/k$ is \emph{$d$-independent} and
determined by the algebraic type of the DSE only.  The dressing
$\gamma_k(d)$ depends on the specific microphysics (via the rank-$k$
bridge~\citep[Section 4]{AmarteifioCFAC2026}) and on the dimension.

\paragraph{Physical content.}
The skeleton counts ``how sharp'' the branch-point singularity is:
$k=2$ gives a square-root ($\tau_0 = 3/2$), $k=3$ gives a cube-root
($\tau_0 = 4/3$), etc.  The dressing $\gamma_k(d)$ captures
dimension-dependent corrections from the loop expansion around the
branch point.

% ==================================================================
\section{The $\Cstr{3}$ slotting claim}
\label{sec:slotting}
% ==================================================================

\begin{slotbox}
\textbf{Claim (slotting).}  The Manna and conserved-DP universality
classes both sit on the $\Cstr{3}$ stratum of
Theorem~\ref{thm:strat}: $k = 3$, skeleton $\tau_0 = 4/3$, with
different $\gamma_3(d=1+1)$ dressings accounting for the observed gap
between their simulation $\tau$ values.
\end{slotbox}

\paragraph{Raw slotting.}
Table~\ref{tab:slotting} lists, for the four best-studied non-DP
universality classes in 1+1d, the published $\tau$, the inferred
$k = 1/(\tau - 1)$, and the nearest integer skeleton stratum.

\begin{table}[h]
\centering
\renewcommand{\arraystretch}{1.2}
\begin{tabular}{lcccc}
\toprule
Class & $\tau$ & $k_{\rm raw}$ & $\Cstr{k}$ & $\tau_0 = 1+1/k$ \\
\midrule
DP 1+1d       & 1.108 & 9.26 & $\Cstr{2}$ & $1.500$ \\
Manna 1+1d    & 1.29  & 3.45 & $\Cstr{3}$ & $1.333$ \\
C-DP 1+1d     & 1.35  & 2.86 & $\Cstr{3}$ & $1.333$ \\
voter 1d      & $\to 1$ & $\to \infty$ & $\Cstr{\infty}$ & $1.000$ \\
\bottomrule
\end{tabular}
\caption{Raw slotting of published $\tau$ values onto integer
  strata.  Values from~\citep{Odor2004,HenkelHinrichsenLubeck}.
  Manna and C-DP both sit within $3\%$ of the $\Cstr{3}$ skeleton.}
\label{tab:slotting}
\end{table}

\paragraph{Signal: Manna and C-DP flank $\Cstr{3}$ from opposite sides.}
The Manna residual $\gamma_3^{\rm Manna} = \tau_{\rm Manna} - \tau_0
\approx -0.04$ and the C-DP residual $\gamma_3^{\rm C-DP} \approx
+0.02$ have \emph{opposite signs} and comparable magnitudes,
consistent with the two classes being a common $\Cstr{3}$ object
differentiated only by their one-loop dimensional dressings.

\paragraph{DP sits on $\Cstr{2}$ with a full 1-loop dressing.}
The raw $k_{\rm raw} = 9.26$ inferred from $\tau_{\rm DP} = 1.108$ is
misleading: DP is an $\Nalg$ DSE with dominant Puiseux order $k=2$
(square-root branch), and the 40\% gap from $\tau_0 = 3/2$ is the
full 1-loop $\gamma_2(d=1)$ at $\eps = 3$.  This is expected, not a
failure of the slotting.

\paragraph{Dyadic / non-dyadic distinction.}
Banderier--Drmota excludes $k=3$ for $\Nalg$ DSEs.  Manna and C-DP
both involve \emph{conserved background fields} coupled to an
absorbing DP order parameter.  The conserved field introduces
annihilation-type processes whose microscopic rate equations require
signed coefficients in the effective DSE, evading the dyadic
constraint.  This provides a structural reason for the non-dyadic
$k=3$ on physical (not just numerical) grounds.

\begin{remark}
The claim is strong: it says Manna and C-DP are \emph{the same
stratum}, differentiated only by dimensional dressing.  The falsifier
is a failure of the predicted $\tau$-gap derivation in
Section~\ref{sec:gamma3}.
\end{remark}

% ==================================================================
\section{The rank-3 bridge and $\gamma_3(d=1+1)$}
\label{sec:gamma3}
% ==================================================================

The dressing $\gamma_k(d)$ in Theorem~\ref{thm:strat} arises from the
rank-$k$ bridge function $B_k$ evaluated at the one-loop DP-like
fixed point and integrated over the $d$-dimensional momentum shell.

\paragraph{Rank-$k$ bridge.}
From~\citep[Prop.~4.3]{AmarteifioCFAC2026},
\begin{equation}\label{eq:bridge_rank_k}
  B_k(d) = \frac{2}{(k-1)!\,(4\pi)^{k/2}}.
\end{equation}
For $k=3$, $d=2$: $B_3(2) = 2/(2! \cdot (4\pi)^{3/2}) = 1/(4\pi)^{3/2}
\approx 0.02247$.

\paragraph{One-loop $\gamma_3$.}
[TODO: insert derivation of $\gamma_3(d)$ from the rank-3 bridge
and the $\Cstr{3}$ one-loop fixed-point amplitude.  The result
should be of the form $\gamma_3(d) = a_3 \eps + O(\eps^2)$ with
$a_3$ a rational multiple of $B_3$ and an integer combinatorial
factor from the cube-root fixed-point profile.]

\paragraph{Prediction.}
The Manna and conserved-DP $\tau$-gap follows from the difference in
their $\gamma_3$ evaluations, which in turn depends only on the
microscopic rate constants of the two processes.  Specifically,
Manna's extra reaction $2A \to \emptyset$ (absorbing pair
annihilation) and C-DP's conservation law enter through the
same one-loop fixed-point amplitude but with different signs.
Writing the predicted gap as
\begin{equation}\label{eq:gap_pred}
  \tau_{\rm C-DP} - \tau_{\rm Manna}
  = \gamma_3^{\rm C-DP}(d{=}2) - \gamma_3^{\rm Manna}(d{=}2),
\end{equation}
the right-hand side is computable without fitting parameters from
Eq.~\eqref{eq:bridge_rank_k} plus the specific one-loop vertex
amplitude for each class.

[TODO: numerical prediction, with target simulation value
$\Delta\tau_{\rm obs} = 1.35 - 1.29 = 0.06$.]

% ==================================================================
\section{Cross-checks against simulations}
\label{sec:sim}
% ==================================================================

Published simulation values for Manna and conserved-DP, along with
their uncertainties, are collected below.  We compare to our
$\gamma_3$-based prediction.

\begin{table}[h]
\centering
\begin{tabular}{lccc}
\toprule
Class & Source & $\tau$ & Error bar \\
\midrule
Manna   & \citep{Odor2004} & 1.29 & $\pm 0.02$ \\
Manna   & \citep{HenkelHinrichsenLubeck} & 1.275 & $\pm 0.01$ \\
C-DP    & \citep{Odor2004} & 1.35 & $\pm 0.03$ \\
C-DP    & \citep{LubeckHucht2002} & 1.338 & $\pm 0.005$ \\
\bottomrule
\end{tabular}
\caption{Published simulation $\tau$ for the two $\Cstr{3}$
  candidates.  The observed gap is
  $\Delta\tau_{\rm obs} \approx 0.06 \pm 0.03$.}
\label{tab:sim}
\end{table}

[TODO: compare prediction from Eq.~\eqref{eq:gap_pred} to the
observed gap $\Delta\tau_{\rm obs}$.  Report relative error.  The
slotting claim is falsifiable: if the prediction is outside the
error bars of the simulations, the claim fails.]

\paragraph{Independent check: DP on $\Cstr{2}$.}
The same framework applied to DP with $k=2$ should give
$\tau_{\rm DP} = 3/2 + \gamma_2(d=1)$, with
$\gamma_2(d=1) \approx -0.39$.  The $\gamma_2(d)$ computation is
standard and has been done at 2~loops in the DP
literature~\citep[Table~4]{Odor2004}; we use it as an independent
check that our $\gamma_k$ machinery is calibrated correctly.

% ==================================================================
\section{Directions we believe do \emph{not} currently slot}
\label{sec:boundary}
% ==================================================================

Scouting four further problems from the CFAC candidate
list~\citep[\texttt{docs/problems.md}]{AmarteifioProblems2026}
returned clear structural reasons why the current stratification
framework does not reach them.  These are not proven no-gos: each
flags a specific extension that would be required.  We state them
here to mark the boundary of the framework honestly.

\subsection{LERW in $d=3$}

Loop-erased random walks in $d=3$ have a scaling exponent
$\nu \approx 1.624$ that has resisted closed form since
Kozma~\citep{Kozma2007} proved it is not equal to $3/2$.
Wilson's algorithm connects LERW to spanning-tree enumeration,
i.e.\ to the Kirchhoff polynomial that CFAC computes natively.
However:
\begin{itemize}[leftmargin=1.2em,itemsep=1pt]
  \item In the thermodynamic limit the LERW generating function
    on $\mathbb{Z}^3$ is a lattice Green's function, not algebraic
    and not $D$-finite; Theorem~\ref{thm:strat} does not apply.
  \item Setting $\tau = \nu$ would require $k = 1/(\nu-1) \approx
    1.60$, neither integer nor dyadic.
\end{itemize}
We believe LERW 3D is not reachable without extending the
stratification framework to non-$D$-finite generating functions.
A legitimate adjacent calculation would be the tropical Kirchhoff
decomposition of Wilson-weight expectations on \emph{finite}
graphs.

\subsection{TAP complexity in spin glasses}

The log-count of critical points of a random spin-glass energy
landscape --- the Bray--Moore / Crisanti--Leuzzi
complexity~\citep{BrayMoore1980,CrisantiLeuzzi2005}, sharpened by
Auffinger--Ben Arous--\v{C}ern\'y~\citep{AuffingerBenArousCerny2013}
and Subag~\citep{Subag2017} for pure $p$-spin --- is controlled by
the Kac--Rice formula $\mathbb{E}|\det(H - f I)|$, where $H$ is the
Hessian, an $N\times N$ GOE matrix.
The CFAC bridge function~\citep[Section 4]{AmarteifioCFAC2026} is
scalar: it handles mass-independent scalar self-energies and rank-$k$
polynomial integrands.  The object that would describe TAP complexity
is a \emph{matrix-valued} bridge --- the operator generalisation
whose scalar limit recovers Eq.~\eqref{eq:bridge_rank_k} but whose
determinantal structure captures the GOE spectral measure.
We believe this is a genuine extension requiring new mathematics
(free probability, $R$-transforms), not a mechanical tensor dressing
of the existing bridge.
TAP complexity is parked pending that extension.

\subsection{KPZ upper-tail universality}

The KPZ equation in 1+1d has an upper-tail rate function of shape
$|h|^{3/2}$ with a coefficient known only for integrable weight
distributions~\citep{CorwinGhosal2020,Krajenbrink}.
Universality beyond integrable weights is one of the top open
problems in integrable probability.
\begin{itemize}[leftmargin=1.2em,itemsep=1pt]
  \item Structural match: $\tau = 3/2$ is exactly the $\Cstr{2}$
    skeleton (square-root branch), the generic Puiseux class of a
    polynomial DSE.  The tilted-DSE machinery in \texttt{tilted.py}
    handles SCGFs of polynomial DSEs.  Reproducing the
    \emph{exponent} is, structurally, cheap.
  \item Missing machinery: the \emph{coefficient} requires a bridge
    function for disorder-averaged random-weight transfer matrices.
    No such bridge is currently implemented; the existing bridges
    cover deterministic Doi--Peliti + MSR only.
\end{itemize}
We believe KPZ upper-tail universality is a ``cautious go'' for the
scoping layer (show that $\tau = 3/2$ comes out of the tilted
directed-polymer DSE at the $\Cstr{2}$ stratum) and a research
programme for the full universality proof.

\subsection{Why listing these matters}
The stratification theorem, as stated, is narrow: it applies to
polynomial DSEs whose generating functions are algebraic or
$D$-finite.  Each negative case above identifies a specific
structural feature (non-$D$-finite GF, matrix-valued bridge,
disorder-averaged bridge) that the current framework does not
handle.  We list them here rather than in a separate note because
the positive result (Manna/C-DP on $\Cstr{3}$) is only
interpretable against the clear boundary of what the framework
currently covers.

% ==================================================================
\section{Discussion}
\label{sec:discussion}
% ==================================================================

If the $\gamma_3$ computation in Section~\ref{sec:gamma3} reproduces
the Manna--C-DP $\tau$-gap within simulation error bars, the
stratification theorem provides a genuinely new organising
principle for non-DP universality: classes are grouped by their
skeleton $k$, with observable $\tau$ differences arising from
calculable dressings rather than being independent phenomenological
parameters.

\paragraph{What is \emph{not} claimed.}
We do not claim to prove the 1+1d Manna or C-DP exponents
analytically.  We claim that the \emph{slotting} --- the
organisation of the universality landscape by $\Cstr{k}$ --- is a
correct coarse structure, falsifiable by the $\gamma_3$ prediction.
A failure of that prediction would rule out the slotting.

\paragraph{Next steps.}
(i) Complete the $\gamma_3(d=1+1)$ calculation in
Section~\ref{sec:gamma3} and confirm or falsify the $\tau$-gap
prediction.
(ii) Extend to the voter class at $k \to \infty$ with log
corrections.
(iii) Test the slotting on a further non-DP class not used in the
hypothesis formation (parity-conserving branching annihilation
random walks, for instance).

% ==================================================================
\begin{thebibliography}{9}

\bibitem[Amarteifio(2026a)]{AmarteifioCFAC2026}
  Amarteifio, S., 2026.
  \textit{CFAC theorem: stratification of Dyson--Schwinger equations
  by Puiseux order}.
  Percolation Labs preprint.

\bibitem[Amarteifio(2026b)]{AmarteifioProblems2026}
  Amarteifio, S., 2026.
  Candidate problems for the CFAC framework.
  \texttt{rdft/docs/problems.md}.

\bibitem[Auffinger, Ben~Arous and \v{C}ern\'y(2013)]{AuffingerBenArousCerny2013}
  Auffinger, A., Ben~Arous, G., and \v{C}ern\'y, J., 2013.
  Random matrices and complexity of spin glasses.
  \textit{Comm. Pure Appl. Math.} \textbf{66}, 165.

\bibitem[Bray and Moore(1980)]{BrayMoore1980}
  Bray, A.~J. and Moore, M.~A., 1980.
  Metastable states in spin glasses.
  \textit{J. Phys. C} \textbf{13}, L469.

\bibitem[Corwin and Ghosal(2020)]{CorwinGhosal2020}
  Corwin, I.\ and Ghosal, P., 2020.
  Lower tail of the KPZ equation.
  \textit{Duke Math. J.} \textbf{169}, 1329.

\bibitem[Crisanti and Leuzzi(2005)]{CrisantiLeuzzi2005}
  Crisanti, A.\ and Leuzzi, L., 2005.
  Spherical 2+$p$ spin-glass model.
  \textit{Phys. Rev. B} \textbf{73}, 014412.

\bibitem[Henkel, Hinrichsen and L\"ubeck]{HenkelHinrichsenLubeck}
  Henkel, M., Hinrichsen, H.\ and L\"ubeck, S., 2008.
  \textit{Non-Equilibrium Phase Transitions, Vol.\ 1:
  Absorbing Phase Transitions}.  Springer.

\bibitem[Krajenbrink(2020)]{Krajenbrink}
  Krajenbrink, A., 2020.
  From Painlev\'e transcendents to the KPZ equation.
  Ph.D. thesis, Ecole Normale Sup\'erieure, Paris.

\bibitem[Kozma(2007)]{Kozma2007}
  Kozma, G., 2007.
  The scaling limit of loop-erased random walk in three dimensions.
  \textit{Acta Math.} \textbf{199}, 29.

\bibitem[L\"ubeck and Hucht(2002)]{LubeckHucht2002}
  L\"ubeck, S.\ and Hucht, A., 2002.
  Absorbing state phase transition in conserved DP.
  \textit{J. Phys. A} \textbf{35}, 4853.

\bibitem[\'Odor(2004)]{Odor2004}
  \'Odor, G., 2004.
  Universality classes in non-equilibrium lattice systems.
  \textit{Rev. Mod. Phys.} \textbf{76}, 663.

\bibitem[Subag(2017)]{Subag2017}
  Subag, E., 2017.
  The complexity of spherical $p$-spin models.
  \textit{Ann. Probab.} \textbf{45}, 3385.

\end{thebibliography}

\end{document}
